% 本文件由 LaTeX 排版 Agent 根据有效 translation.json 自动生成。
% author=John Chrysostom; work_id=1914; title=关于Eutropius的第一篇讲道

\chapter{关于\Person{Eutropius}的第一篇讲道}

% structure: p0001 source=h1 target=omitted reason=page-title-already-rendered-by-parent

\Emph{论太监、贵族与执政官\Person{Eutropius}}\par

1. \BibleQuote{虚而又虚，万事皆虚}——这句话总是合时的，但在现今尤其如此。你执政时那些辉煌的排场如今何在？那些闪烁的火炬何在？那舞蹈和舞者的脚步声、宴席与庆典何在？剧院的华盖与帷幕何在？城中对你的喝彩、竞技场中对你的欢呼、观众的奉承何在？它们逝去了——全部逝去了：一阵风吹在树上，摇落所有叶子，使它全然光秃，连根摇撼；因为那风是如此猛烈，竟震动了树的每一条纤维，大有连根拔起之势。如今你那些虚伪的朋友何在？你的酒宴与晚餐何在？那成群的食客何在？那整日倾泻的酒浆何在？你的厨师所发明的种种美味何在？那些巴结你权位、为讨好你而无所不为的人何在？他们全是黑夜的幻影，随着黎明消失的梦幻；它们是春天的花朵，春天一过便全部凋谢；它们是掠过的阴影——是消散的烟云，是破裂的水泡，是撕裂的蛛网。因此我们不断吟咏这属灵的歌——\BibleQuote{虚而又虚，万事皆虚。}因为这句话应当不断写在我们的墙壁上、衣服上、市场上、家中、街道上、门廊上、入口处，尤其写在每个人的良心上，作为永恒的默想主题。既因为虚假、伪装与表面现象在许多人看来如同真实，所以各人应当每天，无论在晚餐还是早餐时，在社交集会中，对邻人说，并听邻人回答：\BibleQuote{虚而又虚，万事皆虚。}我不是一直在告诉你，财富是逃奴吗？但你却不听。我没有告诉你，它是不感恩的仆人吗？你却不信。现在实际的经历证明了它不仅是逃奴和不感恩的仆人，而且是杀人的凶手，正是它使你如今恐惧战栗。不是曾对你说——你屡次因我说实话而责备我——\Quote{我比那些奉承你的人更爱你？} \Quote{我责备你，比那些讨好你的人更关心你？}我不是还加上这话，说：\BibleQuote{朋友的创伤比敌人的自愿亲吻更可靠。}\BibleRef{箴 27:6}你若接受我的创伤，他们的亲吻就不会使你遭此毁灭；因为我的创伤带来健康，而他们的亲吻却造成了不治之症。如今你的司酒官何在？那些为你开道、在城中对众人不住赞美你的人何在？他们逃走了，否认与你的友谊，借你的患难保全自己。但我却不如此，相反，在你遭难时我不抛弃你，如今你跌倒，我保护并照顾你。你曾视为敌人的教会却敞开怀抱接纳了你；而你曾讨好、并多次因此向我发怒的剧场却出卖并毁灭了你。我从未停止对你说：\Quote{你为什么做这些事？} \Quote{你激怒了教会，正在自取灭亡，}但你却不顾我的一切警告而离去。如今竞技场耗尽了你的财富，磨利了刀剑指向你；而你曾激起其不当义怒的教会却四处奔走，要将你从网罗中拉出。\par

2. 我现在说这些话，并非要践踏一个已倒伏的人，而是希望使那些仍站立的人更加稳固；并非要刺激受伤者的伤口，而是要保护那些尚未受伤的人保持健康；并非要击沉被波涛颠簸的人，而是要教导那些顺风航行的人，使他们不致被淹没。如何能做到这一点呢？通过观察人事的变迁。因为就连这个人，如果他曾畏惧变迁，就不会经历它；但是，既然他自己的良心以及他人的劝告都没有使他改进，那么，你们这些以自己的财富自豪的人，至少要从他的灾祸中获益：因为没有什么比人事更脆弱了。因此，无论用什么词语来表达人事的微不足道，都不足以描述其实际状况——无论是称其为烟、草、梦、春花，还是其他任何名称；它们是如此短暂易逝，比虚无更虚无；而且，除了它们的虚无之外，还带有非常危险的成分，我们眼前就有一个证据。因为谁比这个人更显赫呢？他不是在财富上超越了全世界吗？他不是登上了荣誉的顶峰吗？难道不是所有人都在他面前颤抖畏惧吗？然而，看！他变得比囚犯更悲惨，比奴仆更可怜，比因饥饿而日渐消瘦的乞丐更贫困，每天都要看到磨利的刀剑、罪犯的坟墓，以及刽子手带他走向死亡；他甚至不知道自己是否曾享受过过去的快乐，甚至连阳光也感觉不到，正午时分他的视线也模糊不清，仿佛被最浓密的黑暗所包围。即使我竭尽全力，也无法用言语向你们描述他在时刻等待死亡时必然会经历的痛苦。但实际上，何需我多言呢？他自己已经像一幅可见的图像一样为我们清晰地描绘了这一切。因为昨天，当宫廷的人来找他，想要强行把他拖走时，他逃到祭台那里寻求庇护，那时他的脸色，就像现在一样，比死人好不了多少；他牙齿打颤，全身颤抖，声音颤抖，结结巴巴，总之，他整个人的样子都表明他的灵魂已经僵化了。\par

3. 我现在说这些话，并非要责备他或侮辱他的不幸，而是希望使你们的心对他软化，引导你们怜悯他，并劝你们满足于已经施加的惩罚。因为我们中间有许多不近人情的人，他们可能倾向于责备我允许他进入圣所，我展示他的苦难，是希望通过我的叙述软化他们的铁石心肠。\par

因为请告诉我，亲爱的弟兄，你为什么对我发怒？你说是因为那持续攻击教会的人如今在教会内避难。然而，我们确然应当最大程度地光荣天主，因为他容许他陷入如此巨大的困境，以致既能体验到教会的德能，也能体验到教会的慈爱：她的德能——他因攻击教会的缘故而遭受了这巨大的变故；她的慈爱——他所攻击的教会如今却在他面前举起盾牌，将他接纳在她的翅膀之下，使他完全安全，不记恨任何先前的伤害，而是最慈爱地向他敞开自己的怀抱。因为这事比任何一种凯旋更有荣耀，这是一场辉煌的胜利，它使外邦人和犹太人都感到羞愧，它展示了教会光明的面貌：因为她接纳了她的敌人如同俘虏，却宽恕了他；当所有人都在他孤苦时轻视他时，唯独她像一位慈爱的母亲，用她的外衣遮盖了他，同时对抗君王的愤怒、民众的狂怒以及他们势不可挡的仇恨。这是祭台的装饰。——你说，这是一种奇怪的装饰，当被控告的罪人、敲诈者、盗贼被允许抓住祭台时。不！不要这样说：因为连那个娼妇也曾抓住耶稣的脚，她被最可恶、最不洁的罪所玷污；然而她的行为并不是对耶稣的侮辱，反而增加了对他的钦佩和赞美：因为那不洁的女人并没有伤害那纯洁者，反而那卑贱的娼妇因触摸那纯洁无瑕者而变得纯洁。那么，人哪，不要嫉妒。我们是那被钉者（曾说：\Quote{因为他们不知道他们做的是什么，宽恕他们吧。}）的仆人。 \BibleRef{路 23:34} 但是，你说，他曾通过他的法令和各种法律剥夺了在此避难的权利。是的！然而现在他已从经验中学到了他所作的事，他亲自以他的行为成为第一个违犯法律的人，并且成了全世界的景象；虽然他沉默不语，却从那里向所有人发出警告的声音，说：\Quote{不要做我所做过的事，以免遭受我所受的苦难。}他藉着自己的祸患成为导师，祭台发出极大的光辉，现在因它束缚了这头狮子而激起最大的敬畏；因为任何君王的形象若不仅让人看到君王坐在宝座上，穿着紫袍，戴着王冠，而且还有双手绑在背后的蛮族俯首在地，伏于君王脚下，那就会大大增色。至于没有使用任何具有说服力的言辞，你们自身就是那热忱和人群汇聚的见证。因为今天我们眼前的景象确实辉煌，集会盛大，我看到今天的聚会与圣复活节一样盛大。因此，这人虽未开口，却藉着他的行动发出比号角更清晰的声音，把你们召到这里；今天你们全都聚集于此，少女离开了闺房，主妇离开了内室，男人离开了市场，为的是看见人性被定罪，世俗事物的不稳定被暴露，以及那几天前还容光焕发（因敲诈而得的兴盛就是如此）的娼妇的脸，比任何满脸皱纹的老妇更丑——我说，这脸你可看到它的釉彩和颜料已被逆境如同海绵一样擦去。\par

4. 这便是这场灾难的力量：它使一个曾经显赫出众的人显得最为卑微。若有一个富人进入集会，他从这景象中获益良多：因为当他看见那个曾震撼全世界的人，如今从如此高的权力顶峰被拽下，因恐惧而蜷缩，比兔子或青蛙还害怕，被钉在那边的柱子上，没有锁链，恐惧代替了链条，惊慌失措，颤抖不停，他便收敛傲慢，放下骄矜，获得了关于人事应有的智慧，并带着由榜样所教导的教训离去，这教训正是圣经以训诲所教导的：——\Quote{凡有血气的尽都如草，一切荣华如草上的花：草必枯干，花必凋谢}\BibleRef{依 40:6}，或\Quote{他们必如草快快的枯干，又如青草快快的凋谢}，或\Quote{他的日子如烟消散}，以及所有此类经文。再者，穷人进来看到这景象，并不自卑，也不因自己的贫穷而哀叹，反而感激他的贫穷，因这是他的避难所、平静的港湾和坚固的保障；他看到这些事，常愿留在原位，远胜过暂时拥有所有人的一切，而后自己的生命却陷入危险。你们看见了吗？富人、穷人、尊贵、卑贱、自主的、为奴的，都从这人的避难中获得了不少益处？你们看见了吗？每个人都带着治愈离开，仅凭这景象就得了医治。如何？我是否软化了你们的激情，驱逐了你们的愤怒？我是否熄灭了你们的残忍？我是否使你们变得富有怜悯？我确实认为我已做到；你们的容颜和所流的泪水就是证明。既然你们的顽石已化为深沃的土壤，让我们赶紧结出怜悯的果实，通过俯伏在皇帝面前，或者更确切地说，通过恳求仁慈的天主，来展示丰饶的同情，以软化皇帝的愤怒，使他的心变得温柔，从而赐予我们所求的全部恩惠。因为自从那天这人逃来这里避难起，已经发生了不小的变化；皇帝一得知他匆忙逃入此庇护所，尽管军队在场，因他的恶行而愤怒，要求将他交出处决，皇帝却作了长篇演说，竭力平息士兵的愤怒，声称不仅要考虑他的过犯，也应考虑他可能所做的任何善事，宣布他对后者心怀感激，并准备以同类之谊宽恕他的其他行为。当他们再次催促他报复对皇权的侮辱，呼喊、跳跃、挥舞长矛时，他从温柔的眼睛中流下泪河，提醒他们那人逃往的神圣祭台，最终成功平息了他们的愤怒。\par

5. 此外，让我再补充一些与我们自身相关的论据。倘若皇帝在被侮辱时并不怀恨，而你们这些从未经历过此类事的人却表现出如此愤怒，你们能配得什么宽恕呢？并且，在这次集会解散后，你们将如何处理神圣奥迹，再诵念那命令我们说\Quote{宽恕我们，如同我们也宽恕我们的债务人}\BibleRef{玛 6:12}的祷文，而你们却正要求向你们的债务人报复？他是否对你们造成了极大的伤害和侮辱？我不否认。然而，现在不是审判的时候，而是怜悯的时候；不是要求算账，而是显示慈爱的时候；不是调查诉求，而是允诺的时候；不是判决与报复，而是怜悯与恩宠的时候。因此，愿无人恼怒或烦躁，而让我们一同恳求仁慈的天主，赐予他死亡的缓期，将他从这迫近的毁灭中拯救出来，使他能脱去自己的过犯；让我们同心协力，为教会的缘故，为祭台的缘故，去恳求仁慈的皇帝，将一个人的生命作为奉献给神圣祭台的祭品。我们若如此行，皇帝本人将接纳我们，甚至在得到他的赞许之前，我们已获得天主的认可，祂将因我们的怜悯而赐予丰厚的赏报。因为正如祂拒绝和憎恨残忍与不人道的人，祂也欢迎和喜爱仁慈与仁爱的人；若这样的人是义人，为他编织的冠冕便更加荣耀；若是罪人，祂便赦免他的罪，以此作为他对同伴同情之心的赏报。\Quote{因为祂说：\InnerQuote{我要的是仁慈，不是祭献}}，并且你发现圣经中祂总是在寻求这个，并宣称这是从罪中释放的途径。如此，我们将使祂对我们仁慈，如此我们将释放自己脱离罪过，如此我们将装饰教会，如此我们仁慈的皇帝，正如我已说过的，将称赞我们，全体人民将为我们鼓掌，地极将佩服我们城市的仁爱和温和，世上所有听闻这些事迹的人将颂扬我们。为使我们能享受这些美善，让我们俯伏祈祷恳求，救拔这被掳的、逃难的、祈求者脱离危险，好使我们自己藉我们的主耶稣基督的恩宠和仁慈，获得将来的福乐，愿光荣和权能归于祂，从今世直到永远，世世无尽。阿们。\par

\SourceRecord{Translated by W.R.W. Stephens.；Nicene and Post-Nicene Fathers, First Series；Vol. 9.；Edited by Philip Schaff.；Buffalo, NY: Christian Literature Publishing Co.,；1889.；<http://www.newadvent.org/fathers/1914.htm>.}
